\documentclass{article}
\usepackage{csquotes}

\begin{document}

After almost 20 years of modelling pneumococcal disease, I have started thinking differently about what a ``model'' actually is.

We often imagine the process as:

\begin{quote}
data $\rightarrow$ model $\rightarrow$ prediction $\rightarrow$ intervention $\rightarrow$ new data $\rightarrow$ update.
\end{quote}

But this makes the update look like something that happens \emph{to} a model.

I think it is the other way around.

A functioning model is part of a learning system. And that system exists only because certain processes keep happening:

\begin{itemize}
    \item We observe.
    \item We share information.
    \item We disagree.
    \item We challenge assumptions.
    \item We identify errors.
    \item We correct them.
    \item We update.
\end{itemize}

These are not merely desirable properties of good science.

They are what makes continuous modelling possible in the first place.

If information is withheld, disagreement suppressed, unexpected findings ignored, or mistakes made impossible to acknowledge, the update process doesn't merely become less ethical or less pleasant.

It becomes less capable of learning.

This changes how I think about public health.

We don't only need to build better pneumococcal models. We need to maintain the conditions that allow those models to keep changing when reality changes.

And those conditions are not downstream of the modelling process.

They constitute it.

This also gives me a rather uncomfortable thought about freedom.

We are free to withhold information, ignore evidence or refuse correction. But if we want to remain capable of solving an evolving problem, some of these actions are structurally incompatible with the process we depend on.

And perhaps the same is true of reality.

Our model of reality is never reality itself. It is the temporary product of a process that continuously exposes it to correction.

Perhaps the deepest form of resilience is therefore not having the right model.

It is maintaining the system that allows the model to be wrong.

\end{document}